\chapter{Methodology}

% -------------------------------------------------
\section{System Overview}

The system consists of three primary components:

\begin{enumerate}
    \item \textbf{Financial Terminal Dashboard}: A real-time web interface for traders.
    \item \textbf{Financial Backend Service}: Handles real-world data ingestion, news aggregation, and security.
    \item \textbf{Intelligence Module}: Performs ML-based sentiment analysis on news headlines to generate market sentiment scores.
    \item \textbf{Historical Market Database}: PostgreSQL for persistent storage of market events, news, and sentiment data.
    \item \textbf{Observability Layer}: A real-time monitoring stack using \textbf{Grafana and Prometheus} to provide visibility into the internet protocol stack.
\end{enumerate}

The web client communicates with the backend via secure HTTPS and WebSocket protocols. The system receives high-frequency price ticks from the Yahoo Finance API and live news headlines, processes them for aggregation, persists them to the database, and pushes live updates to the terminal.

% -------------------------------------------------
\section{System Architecture}

The architecture follows a client-server model:

\subsection*{Application Layer (Client Terminal)}
\begin{itemize}
    \item User interface for price and news visualization.
    \item JWT-based secure session management.
\end{itemize}

\subsection*{Processing Layer (Backend)}
\begin{itemize}
    \item Yahoo Finance API Data Ingestion and Normalization.
    \item Financial News Aggregation Engine.
    \item \textbf{ML Sentiment Analyzer} (VADER-based NLP).
    \item WebSocket Real-Time Broadcaster.
\end{itemize}

\subsection*{Storage Layer (Persistent)}
\begin{itemize}
    \item Relational Database (PostgreSQL) for tick and news history.
    \item Structured indexing for rapid historical retrieval.
\end{itemize}

% -------------------------------------------------
\section{Methodology}

The development methodology includes:

\begin{enumerate}
    \item Requirement analysis for real-time financial data structures.
    \item Implementation of the Financial Backend Service with ticker processing logic.
    \item Development of the Yahoo Finance Data Ingestor representing real-world API feeds.
    \item \textbf{Integration of Natural Language Processing (NLP)} for news sentiment analysis.
    \item \textbf{Development of correlation algorithms} to study sentiment-price relationships.
    \item Integration of a relational database for price, news, and sentiment persistence.
    \item Deployment of a WebSocket-based push system for sub-second updates.
    \item Security hardening using TLS and JWT-based authentication.
    \item Performance evaluation of streaming throughput and ML inference latency.
\end{enumerate}

% -------------------------------------------------
\section{Network Architecture and Protocol Stack}

The system incorporates protocols from all four layers of the TCP/IP Internet Protocol Suite to ensure reliable, low-latency financial data delivery.

\subsection{Application Layer}
\begin{itemize}
    \item \textbf{HTTPS (REST APIs)}: Used for fetching news headlines (Finnhub API) and macroeconomic data (FRED API). TLS 1.3 is enforced to ensure data integrity and confidentiality during ingestion.
    \item \textbf{WebSocket (WS)}: Implemented for bidirectional, full-duplex communication between the backend and the trading terminal. This minimizes the overhead of HTTP headers during high-frequency stock price updates.
    \item \textbf{DNS}: Utilized for resolving API provider hostnames (e.g., \texttt{query1.finance.yahoo.com}) into routable IP addresses.
\end{itemize}

\subsection{Transport Layer}
\begin{itemize}
    \item \textbf{TCP}: The fundamental transport protocol for all financial data streams. TCP ensures ordered delivery and retransmission of lost packets—critical for price tickers where missing a tick can invalidate technical indicators.
    \item \textbf{QUIC (Bonus Target)}: Proposed as a future enhancement for the mobile terminal to reduce latency caused by head-of-line blocking and speed up the TLS handshake.
\end{itemize}

\subsection{Network Layer}
\begin{itemize}
    \item \textbf{IPv4/IPv6}: The system uses standard Internet Protocol addressing for cloud communication. Docker's virtual network bridge maps container IP addresses to the host's network stack.
    \item \textbf{ICMP}: Used for network diagnostics and to measure round-trip time (RTT) between the ingestion nodes and financial host providers.
\end{itemize}

\subsection{Data Link / Physical Layer}
\begin{itemize}
    \item \textbf{Ethernet / IEEE 802.11 (Wi-Fi)}: The system is designed to operate over standard local area networks. The Docker environment simulates local hop communication through virtualized network interfaces.
\end{itemize}

% -------------------------------------------------
\section{Tools and Technologies}

The implementation utilizes a professional-grade tech stack focused on networking efficiency and data science performance:

\begin{itemize}
    \item \textbf{Programming Language}: Python 3.11 (Selected for its robust networking libraries).
    \item \textbf{Networking Libraries}: \texttt{requests} (HTTPS), \texttt{websockets} (WS/WSS), \texttt{urllib3}.
    \item \textbf{Data Infrastructure}: PostgreSQL (Persistent storage), Redis (Proposed for real-time caching).
    \item \textbf{Virtualization}: Docker and Docker-Compose for network isolation and environment reproducibility.
    \item \textbf{Monitoring \& Observability}: \textbf{Prometheus} for metrics collection and \textbf{Grafana} for real-time visualization of internet protocol performance.
    \item \textbf{Intelligence Frameworks}: XGBoost for direction prediction, FinBERT (Hugging Face) for NLP.
\end{itemize}

% -------------------------------------------------
\section{Project Timeline (Hypothetical)}

\begin{table}[H]
\centering
\small
\begin{tabular}{|l|p{0.7\textwidth}|}
\hline
\textbf{Week} & \textbf{Milestones and Deliverables} \\ \hline
Week 1 & Requirement Analysis, Internet Protocol selection, and System Architecture design. \\ \hline
Week 2 & Implementation of the Data Ingestion engine (HTTPS/REST) and API integration. \\ \hline
Week 3 & Development of the Real-Time WebSocket server and client-side streaming logic. \\ \hline
Week 4 & Training of the XGBoost prediction models and FinBERT sentiment integration. \\ \hline
Week 5 & Vortex Dashboard development (Streamlit/UI) and Docker containerization. \\ \hline
Week 6 & Performance evaluation, Security hardening, and Final Technical Report. \\ \hline
\end{tabular}
\caption{Hypothetical 6-Week Project Timeline}
\label{tab:timeline}
\end{table}
